% --- Archon physics formalization source begin ---
% archon:physics
% archon:covers PhyXMiniProblems/problem_phyx_mini_0675.lean
% archon:source-report reports/phyx_mini/problem_phyx_mini_0675.source.json

\chapter{Physics problem phyx\_mini\_0675}
\label{ch:PhyXMiniProblems_problem_phyx_mini_0675}

\paragraph{Problem source.}
\#\# Physical scenario

Two thin rods are fastened to the inside of a circular ring as shown in figure. One rod of length \textbackslash{}( D \textbackslash{}) is vertical, and the other of length \textbackslash{}( L \textbackslash{}) makes an angle \textbackslash{}( \textbackslash{}theta \textbackslash{}) with the horizontal. The two rods and the ring lie in a vertical plane. Two small beads are free to slide without friction along the rods.

\#\# Question

Find an expression for the time interval required for the blue bead to slide from point \textbackslash{}textcircled\{B\} to point \textbackslash{}textcircled\{C\} in terms of \textbackslash{}( g \textbackslash{}), \textbackslash{}( L \textbackslash{}), and \textbackslash{}( \textbackslash{}theta \textbackslash{}).

\#\# Answer choices

- A: A: \textbackslash{}( \textbackslash{}sqrt\{\textbackslash{}frac\{2L\}\{g \textbackslash{}tan \textbackslash{}theta\}\} \textbackslash{})

- B: B: \textbackslash{}( \textbackslash{}sqrt\{\textbackslash{}frac\{2L\}\{g \textbackslash{}cos \textbackslash{}theta\}\} \textbackslash{})

- C: C: \textbackslash{}( \textbackslash{}sqrt\{\textbackslash{}frac\{2L\}\{g \textbackslash{}sin \textbackslash{}theta\}\} \textbackslash{})

- D: D: \textbackslash{}( \textbackslash{}sqrt\{\textbackslash{}frac\{L\}\{g \textbackslash{}sin \textbackslash{}theta\}\} \textbackslash{})

\#\# Figure caption (auxiliary; use the image as primary evidence)

The image features a geometric diagram with a circle and three points labeled A, B, and C. 

- The circle is gray with a distinct outline.
- Point A is located at the top of the circle and marked with a red dot. 
- Point B is positioned on the right side of the circle and marked with a blue dot.
- Point C is at the bottom of the circle.
  
Lines and labels:
- A vertical line from point A to point C is marked with the distance \textbackslash{}(D\textbackslash{}).
- A line from point C to point B is marked with the distance \textbackslash{}(L\textbackslash{}).
- A dashed line connects points A and B.
  
An angle \textbackslash{}(\textbackslash{}theta\textbackslash{}) is formed between the horizontal line from point C and line \textbackslash{}(L\textbackslash{}), with \textbackslash{}(\textbackslash{}theta\textbackslash{}) labeled next to the angle.

\#\# Dataset metadata

Dynamics; Physical Model Grounding Reasoning, Multi-Formula Reasoning

\paragraph{Recorded answer/context.}
C: C: \textbackslash{}( \textbackslash{}sqrt\{\textbackslash{}frac\{2L\}\{g \textbackslash{}sin \textbackslash{}theta\}\} \textbackslash{})

\paragraph{Figure/image path.}
/root/proposal\_for\_physic/science-mango/phyx\_mini\_run/phyx\_data/test\_image/675.png

\paragraph{Formalization target.}
create a compiling Lean file with sorry bodies at `PhyXMiniProblems/problem\_phyx\_mini\_0675.lean`. The Lean declarations must preserve the physical quantities, dimensions or dimensional roles, figure labels, governing-law hypotheses, and final relation expressed by this problem.
Use Mathlib/Physlib names found through LeanExplore where available. If a physics API is missing, introduce faithful local abstractions rather than scalar placeholder aliases.

\begin{theorem}[Physics formalization target]
\label{thm:physics:phyx_mini_0675:target}
\lean{PhyXMiniProblems.ProblemPhyXMini0675.problem_phyx_mini_0675}
\uses{def:physics:phyx-mini-0675:phyxminiproblems-problemphyxmini0675-lengthinmeters, def:physics:phyx-mini-0675:phyxminiproblems-problemphyxmini0675-timeinseconds, def:physics:phyx-mini-0675:phyxminiproblems-problemphyxmini0675-accelerationinmeterspersecondsquared, def:physics:phyx-mini-0675:phyxminiproblems-problemphyxmini0675-inclinedrodbeadsetup, def:physics:phyx-mini-0675:phyxminiproblems-problemphyxmini0675-matchesproblemandprimaryfigure, def:physics:phyx-mini-0675:phyxminiproblems-problemphyxmini0675-hasphysicalparameters, def:physics:phyx-mini-0675:phyxminiproblems-problemphyxmini0675-obeysfrictionlessinclinedroddynamics, def:physics:phyx-mini-0675:phyxminiproblems-problemphyxmini0675-obeysconstantaccelerationkinematics}
The assigned autoformalize agent should translate this physics problem into Lean declarations in the covered file, with theorem and lemma proof bodies written as `by sorry`.
\end{theorem}
\begin{proof}
This is an autoformalization task, not a proof task. Produce statements that can later be proved by the physics prover without weakening the physical model.
\end{proof}

% --- Archon declaration topology begin ---
\section{Declaration topology}
\begin{definition}[Length Quantity]
  \label{def:physics:phyx-mini-0675:phyxminiproblems-problemphyxmini0675-lengthquantity}
  \lean{PhyXMiniProblems.ProblemPhyXMini0675.LengthQuantity}
  A nonnegative physical length, independent of the unit used to read it
\end{definition}
\begin{proof}
  This is the declared model definition of Length Quantity.
\end{proof}
\begin{definition}[Time Quantity]
  \label{def:physics:phyx-mini-0675:phyxminiproblems-problemphyxmini0675-timequantity}
  \lean{PhyXMiniProblems.ProblemPhyXMini0675.TimeQuantity}
  A nonnegative physical duration, independent of the unit used to read it
\end{definition}
\begin{proof}
  This is the declared model definition of Time Quantity.
\end{proof}
\begin{definition}[Speed Quantity]
  \label{def:physics:phyx-mini-0675:phyxminiproblems-problemphyxmini0675-speedquantity}
  \lean{PhyXMiniProblems.ProblemPhyXMini0675.SpeedQuantity}
  A signed tangential speed, with dimension length per time
\end{definition}
\begin{proof}
  This is the declared model definition of Speed Quantity.
\end{proof}
\begin{definition}[Acceleration Quantity]
  \label{def:physics:phyx-mini-0675:phyxminiproblems-problemphyxmini0675-accelerationquantity}
  \lean{PhyXMiniProblems.ProblemPhyXMini0675.AccelerationQuantity}
  A signed acceleration, with dimension length per time squared
\end{definition}
\begin{proof}
  This is the declared model definition of Acceleration Quantity.
\end{proof}
\begin{definition}[length In Meters]
  \label{def:physics:phyx-mini-0675:phyxminiproblems-problemphyxmini0675-lengthinmeters}
  \lean{PhyXMiniProblems.ProblemPhyXMini0675.lengthInMeters}
  \uses{def:physics:phyx-mini-0675:phyxminiproblems-problemphyxmini0675-lengthquantity}
  Metre readout of a physical length
\end{definition}
\begin{proof}
  This is the declared model definition of length In Meters.
\end{proof}
\begin{definition}[time In Seconds]
  \label{def:physics:phyx-mini-0675:phyxminiproblems-problemphyxmini0675-timeinseconds}
  \lean{PhyXMiniProblems.ProblemPhyXMini0675.timeInSeconds}
  \uses{def:physics:phyx-mini-0675:phyxminiproblems-problemphyxmini0675-timequantity}
  Second readout of a physical duration
\end{definition}
\begin{proof}
  This is the declared model definition of time In Seconds.
\end{proof}
\begin{definition}[speed In Meters Per Second]
  \label{def:physics:phyx-mini-0675:phyxminiproblems-problemphyxmini0675-speedinmeterspersecond}
  \lean{PhyXMiniProblems.ProblemPhyXMini0675.speedInMetersPerSecond}
  \uses{def:physics:phyx-mini-0675:phyxminiproblems-problemphyxmini0675-speedquantity}
  Metre-per-second readout of a signed speed
\end{definition}
\begin{proof}
  This is the declared model definition of speed In Meters Per Second.
\end{proof}
\begin{definition}[acceleration In Meters Per Second Squared]
  \label{def:physics:phyx-mini-0675:phyxminiproblems-problemphyxmini0675-accelerationinmeterspersecondsquared}
  \lean{PhyXMiniProblems.ProblemPhyXMini0675.accelerationInMetersPerSecondSquared}
  \uses{def:physics:phyx-mini-0675:phyxminiproblems-problemphyxmini0675-accelerationquantity}
  Metre-per-second-squared readout of a signed acceleration
\end{definition}
\begin{proof}
  This is the declared model definition of acceleration In Meters Per Second Squared.
\end{proof}
\begin{definition}[Figure Point]
  \label{def:physics:phyx-mini-0675:phyxminiproblems-problemphyxmini0675-figurepoint}
  \lean{PhyXMiniProblems.ProblemPhyXMini0675.FigurePoint}
  The three labeled points shown on the circular ring
\end{definition}
\begin{proof}
  This is the declared model definition of Figure Point.
\end{proof}
\begin{definition}[Rod]
  \label{def:physics:phyx-mini-0675:phyxminiproblems-problemphyxmini0675-rod}
  \lean{PhyXMiniProblems.ProblemPhyXMini0675.Rod}
  The two rods distinguished by their figure labels
\end{definition}
\begin{proof}
  This is the declared model definition of Rod.
\end{proof}
\begin{definition}[Bead]
  \label{def:physics:phyx-mini-0675:phyxminiproblems-problemphyxmini0675-bead}
  \lean{PhyXMiniProblems.ProblemPhyXMini0675.Bead}
  The two beads distinguished by their colors in the figure
\end{definition}
\begin{proof}
  This is the declared model definition of Bead.
\end{proof}
\begin{definition}[Rod Direction]
  \label{def:physics:phyx-mini-0675:phyxminiproblems-problemphyxmini0675-roddirection}
  \lean{PhyXMiniProblems.ProblemPhyXMini0675.RodDirection}
  Direction category of a rod in the vertical plane
\end{definition}
\begin{proof}
  This is the declared model definition of Rod Direction.
\end{proof}
\begin{definition}[Plane Orientation]
  \label{def:physics:phyx-mini-0675:phyxminiproblems-problemphyxmini0675-planeorientation}
  \lean{PhyXMiniProblems.ProblemPhyXMini0675.PlaneOrientation}
  Orientation of the plane containing a physical component
\end{definition}
\begin{proof}
  This is the declared model definition of Plane Orientation.
\end{proof}
\begin{definition}[Circular Ring]
  \label{def:physics:phyx-mini-0675:phyxminiproblems-problemphyxmini0675-circularring}
  \lean{PhyXMiniProblems.ProblemPhyXMini0675.CircularRing}
  \uses{def:physics:phyx-mini-0675:phyxminiproblems-problemphyxmini0675-lengthquantity}
  A physical circular ring with an otherwise unspecified positive radius
\end{definition}
\begin{proof}
  This is the declared model definition of Circular Ring.
\end{proof}
\begin{definition}[Inclined Rod Bead Setup]
  \label{def:physics:phyx-mini-0675:phyxminiproblems-problemphyxmini0675-inclinedrodbeadsetup}
  \lean{PhyXMiniProblems.ProblemPhyXMini0675.InclinedRodBeadSetup}
  \uses{def:physics:phyx-mini-0675:phyxminiproblems-problemphyxmini0675-lengthquantity, def:physics:phyx-mini-0675:phyxminiproblems-problemphyxmini0675-timequantity, def:physics:phyx-mini-0675:phyxminiproblems-problemphyxmini0675-speedquantity, def:physics:phyx-mini-0675:phyxminiproblems-problemphyxmini0675-accelerationquantity, def:physics:phyx-mini-0675:phyxminiproblems-problemphyxmini0675-figurepoint, def:physics:phyx-mini-0675:phyxminiproblems-problemphyxmini0675-rod, def:physics:phyx-mini-0675:phyxminiproblems-problemphyxmini0675-bead, def:physics:phyx-mini-0675:phyxminiproblems-problemphyxmini0675-roddirection, def:physics:phyx-mini-0675:phyxminiproblems-problemphyxmini0675-planeorientation, def:physics:phyx-mini-0675:phyxminiproblems-problemphyxmini0675-circularring}
  All physical quantities and figure-indexed data for the bead-and-rods setup. constantTangentialAcceleration is the signed acceleration in the direction from each bead's initial point toward its destination. It is data to which the frictionless force law is applied below; its value is not defined to be the requested travel-time expression
\end{definition}
\begin{proof}
  This is the declared model definition of Inclined Rod Bead Setup.
\end{proof}
\begin{definition}[Matches Problem And Primary Figure]
  \label{def:physics:phyx-mini-0675:phyxminiproblems-problemphyxmini0675-matchesproblemandprimaryfigure}
  \lean{PhyXMiniProblems.ProblemPhyXMini0675.MatchesProblemAndPrimaryFigure}
  \uses{def:physics:phyx-mini-0675:phyxminiproblems-problemphyxmini0675-lengthinmeters, def:physics:phyx-mini-0675:phyxminiproblems-problemphyxmini0675-inclinedrodbeadsetup}
  Problem-statement and primary-bitmap information. The vertical diameter-like rod joins A to C and carries label D; the inclined rod joins C to B and carries label L. All three endpoints lie on the ring, both rods are fastened within it, and the dashed AB segment is shown. The red and blue beads begin at A and B, respectively, and both are constrained to travel toward C on their pictured rods
\end{definition}
\begin{proof}
  This is the declared model definition of Matches Problem And Primary Figure.
\end{proof}
\begin{definition}[Has Physical Parameters]
  \label{def:physics:phyx-mini-0675:phyxminiproblems-problemphyxmini0675-hasphysicalparameters}
  \lean{PhyXMiniProblems.ProblemPhyXMini0675.HasPhysicalParameters}
  \uses{def:physics:phyx-mini-0675:phyxminiproblems-problemphyxmini0675-lengthinmeters, def:physics:phyx-mini-0675:phyxminiproblems-problemphyxmini0675-accelerationinmeterspersecondsquared, def:physics:phyx-mini-0675:phyxminiproblems-problemphyxmini0675-inclinedrodbeadsetup}
  Nondegeneracy and the acute physical branch pictured for θ. Gravity and all figure lengths are strictly positive
\end{definition}
\begin{proof}
  This is the declared model definition of Has Physical Parameters.
\end{proof}
\begin{definition}[Obeys Frictionless Inclined Rod Dynamics]
  \label{def:physics:phyx-mini-0675:phyxminiproblems-problemphyxmini0675-obeysfrictionlessinclinedroddynamics}
  \lean{PhyXMiniProblems.ProblemPhyXMini0675.ObeysFrictionlessInclinedRodDynamics}
  \uses{def:physics:phyx-mini-0675:phyxminiproblems-problemphyxmini0675-speedinmeterspersecond, def:physics:phyx-mini-0675:phyxminiproblems-problemphyxmini0675-accelerationinmeterspersecondsquared, def:physics:phyx-mini-0675:phyxminiproblems-problemphyxmini0675-inclinedrodbeadsetup}
  Frictionless constrained dynamics for the blue bead. It is released from rest, and the component of gravity along rod BC is g sin θ, constant toward C because the rod is straight
\end{definition}
\begin{proof}
  This is the declared model definition of Obeys Frictionless Inclined Rod Dynamics.
\end{proof}
\begin{definition}[Obeys Constant Acceleration Kinematics]
  \label{def:physics:phyx-mini-0675:phyxminiproblems-problemphyxmini0675-obeysconstantaccelerationkinematics}
  \lean{PhyXMiniProblems.ProblemPhyXMini0675.ObeysConstantAccelerationKinematics}
  \uses{def:physics:phyx-mini-0675:phyxminiproblems-problemphyxmini0675-lengthinmeters, def:physics:phyx-mini-0675:phyxminiproblems-problemphyxmini0675-timeinseconds, def:physics:phyx-mini-0675:phyxminiproblems-problemphyxmini0675-speedinmeterspersecond, def:physics:phyx-mini-0675:phyxminiproblems-problemphyxmini0675-accelerationinmeterspersecondsquared, def:physics:phyx-mini-0675:phyxminiproblems-problemphyxmini0675-inclinedrodbeadsetup}
  The constant-acceleration displacement law L = u t + (1/2) a t² for the blue bead's complete trip from B to C
\end{definition}
\begin{proof}
  This is the declared model definition of Obeys Constant Acceleration Kinematics.
\end{proof}
\begin{lemma}[blue Bead Travel Time squared]
  \label{lem:physics:phyx-mini-0675:phyxminiproblems-problemphyxmini0675-bluebeadtraveltime-squared}
  \lean{PhyXMiniProblems.ProblemPhyXMini0675.blueBeadTravelTime_squared}
  \uses{def:physics:phyx-mini-0675:phyxminiproblems-problemphyxmini0675-lengthinmeters, def:physics:phyx-mini-0675:phyxminiproblems-problemphyxmini0675-timeinseconds, def:physics:phyx-mini-0675:phyxminiproblems-problemphyxmini0675-accelerationinmeterspersecondsquared, def:physics:phyx-mini-0675:phyxminiproblems-problemphyxmini0675-inclinedrodbeadsetup, def:physics:phyx-mini-0675:phyxminiproblems-problemphyxmini0675-hasphysicalparameters, def:physics:phyx-mini-0675:phyxminiproblems-problemphyxmini0675-obeysfrictionlessinclinedroddynamics, def:physics:phyx-mini-0675:phyxminiproblems-problemphyxmini0675-obeysconstantaccelerationkinematics}
  The governing laws first determine the square of the blue bead's travel time. This is a derived intermediate result, not an assumption field
\end{lemma}
\begin{proof}
  The conclusion follows from the governing relations and figure readouts named in the statement of blue Bead Travel Time squared.
\end{proof}
% --- Archon declaration topology end ---
% --- Archon physics formalization source end ---
